% Part: first-order-logic
% Chapter: proof-systems
% Section: tableaux

\documentclass[../../../include/open-logic-section]{subfiles}

\begin{document}

\iftag{FOL}
      {\olfileid{fol}{prf}{tab}}
      {\olfileid{pl}{prf}{tab}}

\olsection{\usetoken{P}{tableau}}

بھاویں بہت سارے !!{derivation} نظام !!{sentence}s دے بندوبستاں نال کم
کردے نیں، !!{tableau}s !!{signed formula}s نال کم کردے نیں۔
!!^a{signed formula} اک جوڑی اے جیہڑی صداقتی قدر دی اک علامت
($\True$ یا $\False$) تے !!a{sentence} اُتے مشتمل ہُندی اے
\[
\sFmla{\True}{!A} \text{ یا } \sFmla{\False}{!A}.
\]
!!^a{tableau} وچ !!{signed formula}s نوں ہِٹھاں ول شاخاں کڈھن والے
درخت دے روپ وچ ترتیب دتا جاندا اے۔ ایہ کجھ \emph{مفروضیاں} نال شروع
ہُندا اے تے اوہناں !!{signed formula}s نال اگے ودھدا اے جیہڑے استنباط
دے قاعدیاں وچوں کوئی اک لاگو کر کے اپنے توں اُتّے والے
!!{signed formula}s وچوں کسے اک توں حاصل ہُندے نیں۔ ہر قاعدہ سانوں کسے شاخ دے سرے
اُتے اک یا ودھ !!{signed formula}s، یا نال نال دو !!{signed formula}s
شامل کرن دی اجازت دیندا اے---ایس دوجی صورت وچ شاخ دو حصیاں وچ ونڈی
جاندی اے تے شامل کیتے دو !!{signed
  formula}s اوہناں دو شاخاں دے سرے بن جاندے نیں۔

کسی مرکب !!{signed formula} اُتے لاگو کیتا قاعدہ اوہ !!{signed formula}s
شامل کردا اے جیہڑے عین ذیلی !!{formula}s ہُندے نیں۔ ایہ
قاعدے جوڑیاں وچ آؤندے نیں، دوہاں علامتاں وچوں ہر اک لئی اک قاعدہ۔
مثال لئی، $\TRule{\True}{\land}$ دا قاعدہ
$\sFmla{\True}{!A \land !B}$ اُتے لاگو ہُندا اے، تے ایہ
دونیں !!{signed formula}s $\sFmla{\True}{!A}$ تے~$\sFmla{\True}{!B}$
نوں $\sFmla{\True}{!A \land !B}$ والی کسے وی شاخ دے سرے اُتے شامل
کرن دی اجازت دیندا اے؛ تے $\TRule{\False}{!A \land !B}$ دا قاعدہ
$\sFmla{\False}{!A}$ تے $\sFmla{\False}{!B}$ نوں نال نال شامل کر کے
شاخ نوں دو حصیاں وچ ونڈن دی اجازت دیندا اے۔ !!^a{tableau} بند اے جے
اوہدی ہر اک شاخ وچ !!{signed formula}s دی اک میل کھاندی جوڑی
$\sFmla{\True}{!A}$ تے $\sFmla{\False}{!A}$ شامل ہووے۔

!!{tableau}s اُتے مبنی $\Proves$ تعلق دی تعریف ایویں کیتی جاندی اے:
$\Gamma \Proves !A$ اُدوں تے صرف اُدوں جدوں کوئی محدود سیٹ~$\Gamma_0 = \{!B_1,
\dots, !B_n\} \subseteq \Gamma$ ایہو جہا ہووے کہ مفروضیاں
\[
\{\sFmla{\False}{!A}, \sFmla{\True}{!B_1}, \dots, \sFmla{\True}{!B_n}\}
\]
لئی کوئی بند !!{tableau} ہووے۔ مثال لئی، ایہ بند !!{tableau} وکھاؤندا
اے کہ $\Proves (!A \land !B) \lif !A$:
\begin{oltableau}
  [\sFmla{\False}{(\formula{A} \land \formula{B}) \lif \formula{A}}, just = \TAss
    [\sFmla{\True}{\formula{A} \land \formula{B}}, just = {\TRule{\False}{\lif}[1]}
      [\sFmla{\False}{\formula{A}}, just = {\TRule{\False}{\lif}[1]}
        [\sFmla{\True}{\formula{A}}, just = {\TRule{\True}{\lif}[2]}
          [\sFmla{\True}{\formula{B}}, just = {\TRule{\True}{\lif}[2]}, close]
        ]
      ]
    ]
  ]
\end{oltableau}

!!{tableau} حساب وچ سیٹ $\Gamma$ متضاد اے جے مفروضیاں
\[
\{\sFmla{\True}{!B_1}, \dots, \sFmla{\True}{!B_n}\}
\]
لئی، کسے $!B_i \in \Gamma$ واسطے، کوئی بند !!{tableau} ہووے۔

!!^{tableau}s نوں 1950 دی دہائی وچ Evert Beth تے Jaakko Hintikka
نے اک دوجے توں آزاد ہو کے ایجاد کیتا، تے Raymond Smullyan نے اوہناں
نوں سادہ بنا کے عام کیتا۔ ایہ ورتن وچ بہت سَوکھے نیں، کیوں جو
!!a{tableau} بنانا بہت منظم کارروائی اے۔ !!{tableau}s دی منظم نوعیت
پاروں ایہناں نوں کمپیوٹر راہیں لاگو کرنا وی سَوکھا اے۔ پر
!!a{tableau} اکثر پڑھن وچ اوکھا ہُندا اے تے اوہناں دا ثبوتاں نال تعلق
ویکھنا کدی کدی سَوکھا نہیں ہُندا۔ ایہ طریقہ کافی عام وی اے تے بہت
ساریاں منطقاں دے !!{tableau} نظام نیں۔ !!^{tableau}s سانوں اوہ
!!{structure}s لبھن وچ وی مدد دیندے نیں جیہڑے دتے ہوئے !!{sentence}s
(یا اوہناں دے سیٹاں) نوں پورا کردے نیں: جے سیٹ پورا ہون جوگ اے تاں
اوہدے لئی کوئی بند !!{tableau} نہیں ہووے گا، یعنی ہر !!{tableau} وچ
اک کھلی شاخ ہووے گی۔ پورا کرن والی !!{structure} نوں کسے کھلی شاخ
توں ``پڑھ کے کڈھیا'' جا سکدا اے، بشرطیکہ اوس شاخ اُتے ہر اوہ قاعدہ
لاگو کیتا جا چکیا ہووے جیہڑا لاگو ہو سکدا اے۔ سیکوئنٹ حساب نال وی
ایس دا بہت نیڑا تعلق اے: بنیادی طور تے، بند !!{tableau} سیکوئنٹ حساب
دا اک مختصر کیتا !!{derivation} اے جیہڑا اُلٹا لکھیا ہُندا اے۔

\end{document}
